\section{Tied orders, empirical laws, and conjectures}
\label{sec:conjectures}

\subsection{Tied maximizers}

The extremal behavior at orders congruent to two modulo four is structurally
different from the uniform double-ladder recurrence.

\begin{computationalresult}[Certified tied-order structure]
\label{res:tied-structure}
Every one of the 17 tied maximizers at orders (22,26,30,34) contains a
3-cycle component in the face-adjacency graph induced by quadrilateral faces,
as well as at least one strict ladder component.  The full quadrilateral
component signatures are pairwise distinct within each order and reproduce the
tie counts (2,5,3,7).

Twenty-five exact square-gadget certificates link these graphs to earlier
examples.  Thirteen links are literal ladder-cell insertions.  In the other
twelve, the same four-vertex gadget is inserted on two suitable disjoint edges
of a 6-, 8-, or 10-face end-cap region.  Thus these graphs are not asserted to
be members of one \(D(a,b)\) recurrence.
\end{computationalresult}

\begin{figure}[htbp]
\centering
\begin{tabular}{cc}
\includegraphics[width=.44\textwidth]{tied22a.png} &
\includegraphics[width=.44\textwidth]{tied22b.png}
\end{tabular}
\caption{The two certified order-22 maximizers.  Purple marks the
quadrilateral 3-cycle component; blue and red distinguish ladder rails and
rungs.  Their different component signatures illustrate why the tied branch is
not identified with a single double-ladder skeleton.}
\label{fig:tied-order-22}
\end{figure}

\subsection{Empirical hsep formula}

The six independently certified values from orders 16 through 36 agree with
\begin{equation}
\label{eq:empirical-hsep}
 \hsep(D(a,b))=\frac{ab+2a+2b+3}{2}.
\end{equation}
The supporting parameter pairs and values are
\[
\begin{array}{c|rrrrrr}
(a,b)&(3,3)&(5,3)&(5,5)&(7,5)&(7,7)&(9,7)\\ \hline
\hsep&12&17&24&31&40&49.
\end{array}
\]
The formula does not match the two degenerate initial members \(D(1,1)\) and
\(D(3,1)\).

\begin{conjecture}[Double-ladder hsep formula]
\label{conj:hsep-formula}
For odd \(a,b\ge3\), equation \cref{eq:empirical-hsep} holds.
\end{conjecture}

This conjecture is not a consequence of the finite expansion certificates.
In the five primary packing certificates at orders 20 through 36, exactly
\(k-2\) packing requirements have both edges internal to the ladders, while the
remaining two have only their included edge internal.  No invariant placement
or certificate-lifting theorem has yet been found.

\begin{conjecture}[Extremal continuation]
Balanced or nearly balanced double ladders continue to be extremal at larger
orders divisible by four, and eventually are uniquely extremal there.
\end{conjecture}

Even if \cref{conj:hsep-formula} is true, it does not prove this extremal
statement; all competing Barnette graphs at a given order would still need an
upper certificate.

\begin{computationalresult}[Prospective double-ladder test]
\label{res:prospective-ladders}
Before constructing or analyzing any new graph, the formulas predicted
\[
\begin{array}{c|rr}
G & |\Hcal(G)| & \hsep(G)\\ \hline
D(9,9) & 86 & 60\\
D(11,9) & 104 & 71\\
D(11,11) & 126 & 84.
\end{array}
\]
The prediction lock was then left unchanged.  Each graph was constructed from
the certified (D(9,7)) member by successive four-vertex square insertions.
For all three graphs, a path-based enumeration and an independent exhaustive
perfect-matching enumeration produced the predicted complete Hamiltonian
universe.  Each predicted \(\hsep\) value has an explicit primal cover and a
matching packing of the same size, checked by a standard-library verifier.
\end{computationalresult}

Thus the two family formulas survive three genuinely prospective tests, but a
finite continuation does not prove \cref{conj:hsep-formula}.  No order-40,
order-44, or order-48 census or extremal certification was performed.  In
particular, this manuscript does not claim values for \(\MB(40)\),
\(\MB(44)\), or \(\MB(48)\).  A separate linear fit to only
the three tied-order values at (26,30,34) predicts \(\MB(38)=39\); because the
tied expansion graph branches and merges, that extrapolation is considered
especially weak.  No order-38 computation is reported here.
